\chapter{Aires et volumes}\label{ch:g7:areas}

L'année 6 a mesuré rectangles, \omterm{def:g6:shapes:triangles}{triangles rectangles} et \omterm{def:g5:solids:def}{pavés}
(\cref{ch:g6:measure}). Avec une paire de ciseaux --- couper un
morceau ici, coller là --- les mêmes idées donnent maintenant les
\omterm{def:g6:measure:area}{aires} des \omterm{def:g7:central:parallelogram}{parallélogrammes}, de \emph{tous} les triangles et du disque,
et les \omterm{def:g6:measure:volume}{volumes} des \omterm{def:g7:areas:prism}{prismes} et des \omterm{def:g7:areas:prism}{cylindres}.

\section{Aire d'un parallélogramme}

\begin{theorem}[Parallélogramme]\label{thm:g7:areas:parallelogram}
Un \omterm{def:g7:central:parallelogram}{parallélogramme} de côté de longueur $b$ (sa \emph{base}) et de
distance $h$ entre ce côté et le côté opposé (sa \emph{hauteur}) a
pour \omterm{def:g6:measure:area}{aire}
\[
A = b \times h .
\]
\end{theorem}

\begin{proof}[Preuve par les ciseaux]
Couper le \omterm{def:g6:shapes:triangles}{triangle rectangle} qui dépasse d'un côté et le coller de
l'autre~: le \omterm{def:g7:central:parallelogram}{parallélogramme} devient un rectangle $b \times h$ de
même \omterm{def:g6:measure:area}{aire}.
\end{proof}

\begin{omfigure}
\begin{tikzpicture}[scale=0.82]
  \fill[omDef!14] (0,0) -- (4,0) -- (5.2,1.8) -- (1.2,1.8) -- cycle;
  \draw[thick] (0,0) -- (4,0) -- (5.2,1.8) -- (1.2,1.8) -- cycle;
  \fill[omThm!25] (0,0) -- (1.2,1.8) -- (1.2,0) -- cycle;
  \draw[omThm, thick] (0,0) -- (1.2,1.8) -- (1.2,0) -- cycle;
  \draw[dashed] (1.2,0) -- (1.2,1.8);
  \node[below, font=\small] at (2.0,0) {$b$};
  \draw[Stealth-Stealth] (6.0,0) -- (6.0,1.8)
    node[midway, right, font=\small] {$h$};
  \begin{scope}[xshift=8.2cm]
  \fill[omDef!14] (0,0) rectangle (4,1.8);
  \draw[thick] (0,0) rectangle (4,1.8);
  \fill[omThm!25] (2.8,0) -- (4,1.8) -- (4,0) -- cycle;
  \draw[omThm, thick] (2.8,0) -- (4,1.8) -- (4,0) -- cycle;
  \node[below, font=\small] at (2.0,0) {$b$};
  \end{scope}
\end{tikzpicture}

{\small La preuve par les ciseaux~: glisser le triangle rouge de
l'extrémité gauche à l'extrémité droite transforme le \omterm{def:g7:central:parallelogram}{parallélogramme}
en rectangle --- même \omterm{def:g6:measure:area}{aire}, $b \times h$.}
\end{omfigure}

\begin{remark}
La hauteur est la distance \emph{\omterm{def:g6:lines:perp}{perpendiculaire}} entre les deux
bases, et non la longueur du côté oblique~! Un \omterm{def:g7:central:parallelogram}{parallélogramme} très
oblique peut avoir de longs côtés et une petite \omterm{def:g6:measure:area}{aire}.
\end{remark}

\section{Aire d'un triangle}

\begin{theorem}[Triangle]\label{thm:g7:areas:triangle}
Un triangle de base $b$ et de hauteur correspondante $h$ (la distance
\omterm{def:g6:lines:perp}{perpendiculaire} du \omterm{def:g5:solids:def}{sommet} opposé à la droite-support de la base) a
pour \omterm{def:g6:measure:area}{aire}
\[
A = \frac{b \times h}{2} .
\]
\end{theorem}

\begin{proof}
Deux copies du triangle, l'une tournée d'un demi-tour autour du
milieu d'un côté, s'assemblent en un \omterm{def:g7:central:parallelogram}{parallélogramme} de base $b$ et de
hauteur $h$ (c'est la \omterm{def:g7:central:def}{symétrie centrale} à l'œuvre,
\cref{prop:g7:central:preserved}). Le triangle en est la \omterm{def:g3:division:half}{moitié}~:
$\frac{bh}{2}$.
\end{proof}

\begin{omfigure}
\begin{tikzpicture}[scale=0.9]
  \fill[omDef!14] (0,0) -- (4,0) -- (1.4,2.1) -- cycle;
  \draw[thick] (0,0) -- (4,0) -- (1.4,2.1) -- cycle;
  \fill[omThm!12] (4,0) -- (1.4,2.1) -- (5.4,2.1) -- cycle;
  \draw[thick, omThm, dashed] (4,0) -- (5.4,2.1) -- (1.4,2.1);
  \draw[dashed] (1.4,0) -- (1.4,2.1);
  \draw[thin] (1.4,0) ++(-0.2,0) -- ++(0,0.2) -- ++(0.2,0);
  \node[below, font=\small] at (2,0) {$b$};
  \node[right, font=\small] at (1.42,1.0) {$h$};
\end{tikzpicture}

{\small Un triangle et sa copie tournée d'un demi-tour (en pointillés)
forment un \omterm{def:g7:central:parallelogram}{parallélogramme}~: l'\omterm{def:g6:measure:area}{aire} du triangle est
$\frac{b\times h}{2}$. N'importe lequel des trois côtés peut servir
de base --- avec \emph{sa propre} hauteur.}
\end{omfigure}

\begin{example}\label{ex:g7:areas:triangle}
Un triangle a une base de $9$ cm et la hauteur correspondante mesure
$4$ cm~:
\[
A = \frac{9 \times 4}{2} = \frac{36}{2} = 18 \text{ cm}^2 .
\]
Pour un triangle \emph{rectangle}, les deux côtés de l'\omterm{def:g3:shapes:rightangle}{angle droit}
sont une base et sa hauteur~: on retrouve la formule de
\cref{prop:g6:measure:formulas}.
\end{example}

\section{Aire d'un disque}

\begin{theorem}[Disque]\label{thm:g7:areas:disk}
Un disque de \omterm{def:g3:shapes:shapes}{rayon} $r$ a pour \omterm{def:g6:measure:area}{aire}
\[
A = \pi r^2 .
\]
\end{theorem}

\begin{proof}[Idée de preuve]
Découper le disque en de nombreuses parts minces égales et les
disposer tête-bêche~: elles forment presque un \omterm{def:g7:central:parallelogram}{parallélogramme} de
hauteur $r$ dont la base est la \omterm{def:g3:division:half}{moitié} de la circonférence, $\pi r$
(\cref{prop:g6:measure:circle}). Son \omterm{def:g6:measure:area}{aire} est proche de
$\pi r \times r = \pi r^2$, et l'approximation devient parfaite
lorsque les parts s'amincissent. Une preuve complète demande le calcul
intégral, dans le volume de lycée.
\end{proof}

\begin{omfigure}
\begin{tikzpicture}[scale=0.85]
  \foreach \a in {0,45,...,315}
    \fill[omDef!20, draw=omDef] (0,0) -- (\a:1.4)
      arc (\a:\a+45:1.4) -- cycle;
  \begin{scope}[xshift=4.3cm, yshift=-0.15cm]
  \foreach \i in {0,1,2,3} {
    \fill[omDef!20, draw=omDef]
      ({\i*1.1},0) -- ({\i*1.1+0.55},1.4) -- ({\i*1.1+1.1},0) -- cycle;
    \fill[omDef!20, draw=omDef]
      ({\i*1.1+0.55},1.4) -- ({\i*1.1+1.1},0) -- ({\i*1.1+1.65},1.4)
      -- cycle;
  }
  \draw[Stealth-Stealth] (5.2,0) -- (5.2,1.4)
    node[midway, right, font=\small] {$\approx r$};
  \node[below, font=\small] at (2.5,-0.1) {base $\approx \pi r$};
  \end{scope}
\end{tikzpicture}

{\small Trancher le disque et dérouler les parts~: presque un
\omterm{def:g7:central:parallelogram}{parallélogramme} de base $\pi r$ (la \omterm{def:g3:division:half}{moitié} du bord) et de hauteur
$r$.}
\end{omfigure}

\begin{example}\label{ex:g7:areas:disk}
Un plateau de table ronde a pour \omterm{def:g3:shapes:shapes}{rayon} $60$ cm. Son \omterm{def:g6:measure:area}{aire} est
$\pi \times 60^2 = 3600\pi \approx 11\,310$ cm$^2$, un peu plus de
$1.1$ m$^2$. Attention à la \omterm{ex:g1:subtraction:difference}{différence}~: le \emph{\omterm{def:g6:measure:perimeter}{périmètre}}
$2\pi r$ utilise $r$, l'\emph{\omterm{def:g6:measure:area}{aire}} $\pi r^2$ utilise $r^2$.
\end{example}

\section{Prismes et cylindres}

\begin{definition}[Prisme, cylindre]\label{def:g7:areas:prism}
Un \emph{prisme}\index{prisme} est un \omterm{def:g5:solids:def}{solide} à deux \omterm{def:g5:solids:def}{faces} polygonales
identiques et \omterm{def:g6:lines:perp}{parallèles} (les \emph{bases}) jointes par des
rectangles~; un \emph{cylindre}\index{cylindre} est la même chose avec
des disques pour bases. La \emph{hauteur} est la distance entre les
deux bases.
\end{definition}

\begin{theorem}[Volumes]\label{thm:g7:areas:volume}
Pour un \omterm{def:g7:areas:prism}{prisme} ou un \omterm{def:g7:areas:prism}{cylindre} d'\omterm{def:g6:measure:area}{aire} de base $B$ et de hauteur $h$~:
\[
V = B \times h .
\]
\end{theorem}
\admitted

\begin{example}\label{ex:g7:areas:volume}
Une tente est un \omterm{def:g7:areas:prism}{prisme} couché sur le côté~: ses bases sont des
triangles de base $2$ m et de hauteur $1.5$ m, et la tente mesure
$3$ m de long.
\begin{enumerate}
  \item \omterm{def:g6:measure:area}{Aire} de base~: $B = \frac{2 \times 1.5}{2} = 1.5$ m$^2$.
  \item \omterm{def:g6:measure:volume}{Volume}~: $V = B \times h = 1.5 \times 3 = 4.5$ m$^3$.
\end{enumerate}
Une boîte de conserve de \omterm{def:g3:shapes:shapes}{rayon} $4$ cm et de hauteur $11$ cm~:
$V = \pi \times 4^2 \times 11 = 176\pi \approx 553$ cm$^3$, à peu
près un demi-litre.
\end{example}

\section{Exercices}

\begin{exercise}[$\star$]\label{exo:g7:areas:1}
Calculer l'\omterm{def:g6:measure:area}{aire} d'un \omterm{def:g7:central:parallelogram}{parallélogramme} de base $8$ cm et de hauteur
$5$ cm~; d'un autre de base $6.5$ m et de hauteur $4$ m.
\end{exercise}

\begin{exercise}[$\star$]\label{exo:g7:areas:2}
Un \omterm{def:g7:central:parallelogram}{parallélogramme} a des côtés de $10$ cm et $6$ cm, et la hauteur
relative à la base de $10$ cm mesure $4$ cm. Calculer son \omterm{def:g6:measure:area}{aire}.
Pourquoi la réponse n'est-elle pas $60$ cm$^2$~?
\end{exercise}

\begin{exercise}[$\star$]\label{exo:g7:areas:3}
Calculer l'\omterm{def:g6:measure:area}{aire} d'un triangle de base $12$ cm et de hauteur $7$ cm~;
d'un \omterm{def:g6:shapes:triangles}{triangle rectangle} de côtés de l'\omterm{def:g3:shapes:rightangle}{angle droit} $9$ et $6$~; d'un
triangle de base $5.5$ m et de hauteur $2$ m.
\end{exercise}

\begin{exercise}[$\star$]\label{exo:g7:areas:4}
Calculer l'\omterm{def:g6:measure:area}{aire} d'un disque de \omterm{def:g3:shapes:shapes}{rayon} $5$ cm, puis d'un demi-disque de
\omterm{def:g3:shapes:shapes}{rayon} $10$ cm (valeurs exactes avec $\pi$, puis \omterm{def:g6:decimals:rounding}{arrondies} à l'unité).
\end{exercise}

\begin{exercise}[$\star$]\label{exo:g7:areas:5}
Calculer le \omterm{def:g6:measure:perimeter}{périmètre} \emph{et} l'\omterm{def:g6:measure:area}{aire} d'un disque de \omterm{def:g3:shapes:shapes}{rayon} $3$ cm.
Laquelle des deux réponses est en cm et laquelle en cm$^2$~?
\end{exercise}

\begin{exercise}[$\star$]\label{exo:g7:areas:6}
Calculer le \omterm{def:g6:measure:volume}{volume} d'un \omterm{def:g7:areas:prism}{prisme} d'\omterm{def:g6:measure:area}{aire} de base $24$ cm$^2$ et de
hauteur $10$ cm~; d'un \omterm{def:g7:areas:prism}{cylindre} de \omterm{def:g3:shapes:shapes}{rayon} $3$ cm et de hauteur $8$ cm
(exact, puis \omterm{def:g6:decimals:rounding}{arrondi}).
\end{exercise}

\begin{exercise}[$\star$]\label{exo:g7:areas:7}
Un triangle a pour \omterm{def:g6:measure:area}{aire} $28$ cm$^2$ et pour base $8$ cm. Trouver la
hauteur correspondante, en écrivant d'abord l'équation.
\end{exercise}

\begin{exercise}[$\star\star$]\label{exo:g7:areas:8}
Tracer un triangle quelconque et mesurer soigneusement les trois
paires base--hauteur. Calculer $\frac{b \times h}{2}$ pour chaque
paire~: les trois résultats devraient coïncider (à l'erreur de
mesure près). Pourquoi~?
\end{exercise}

\begin{exercise}[$\star\star$]\label{exo:g7:areas:9}
Un jardin en forme de trapèze peut se découper, par une diagonale, en
deux triangles partageant la même hauteur $h = 20$ m, de bases $30$ m
et $18$ m. Calculer son \omterm{def:g6:measure:area}{aire} totale. Peut-on deviner une formule
générale pour le trapèze~?
\end{exercise}

\begin{exercise}[$\star\star$]\label{exo:g7:areas:10}
Un vase cylindrique de \omterm{def:g3:shapes:shapes}{rayon} $6$ cm contient de l'eau jusqu'à une
hauteur de $15$ cm. Toute l'eau est versée dans une boîte vide de
base $18$ cm $\times$ $12$ cm. Quelle hauteur d'eau atteint-on dans
la boîte~? (Même \omterm{def:g6:measure:volume}{volume}, base différente.)
\end{exercise}

\begin{exercise}[$\star\star$]\label{exo:g7:areas:11}
Une pizza de \omterm{def:g4:geometry:circle}{diamètre} $30$ cm coûte $9$ euros~; une de \omterm{def:g4:geometry:circle}{diamètre}
$40$ cm coûte $14$ euros. Calculer les deux \omterm{def:g6:measure:area}{aires} et le prix pour
$100$ cm$^2$. Quelle pizza est la meilleure affaire~?
\end{exercise}

\begin{exercise}[$\star\star\star$]\label{exo:g7:areas:12}
Les deux côtés de l'\omterm{def:g3:shapes:rightangle}{angle droit} d'un \omterm{def:g6:shapes:triangles}{triangle rectangle} mesurent
$6$ cm et $8$ cm, et son hypoténuse mesure $10$ cm. Calculer son \omterm{def:g6:measure:area}{aire}
avec les côtés de l'\omterm{def:g3:shapes:rightangle}{angle droit}, puis utiliser cette \omterm{def:g6:measure:area}{aire} pour
trouver la hauteur relative à l'\emph{hypoténuse}.
\end{exercise}
